\cleardoublepage

\section{结论与展望}

\subsection{本文总结}
本文针对具身智能中的视觉语言导航（VLN）任务，提出并实现了一套结合多模态推理能力与强化策略优化的导航系统——VLN-R1。

本文的主要研究成果与贡献总结如下：
\begin{enumerate}
    \item **特征适配与感知增强**：通过引入 DINOv2 视觉编码器，利用其强大的几何空间感知能力，解决了传统 CLIP 模型在连续环境定位精度不足的问题。同时，设计了 4-View 视场对齐策略，成功弥补了仿真环境与实车硬件之间的 FOV 差异。
    \item **显式思维链推理范式**：通过自动化数据引擎 VLN-Ego，为模型赋予了类似人类的“先思考、后动作”的推理能力。这种范式不仅提升了复杂指令下的导航导航成功率，也极大地增强了系统的可解释性。
    \item **强化策略优化算法的创新应用**：首次将 GRPO 算法引入 VLN 领域。相比传统的 RL 方法，GRPO 无需显式的价值网络，显著降低了训练方差并节省了算力资源。通过多维奖励函数的约束，模型在逻辑一致性上表现出更强的鲁棒性。
    \item **实地闭环部署与验证**：在 TurtleBot4 机器人平台上完成了“端-边-云”分布式部署。实验证明，VLN-R1 在真实办公环境下能够有效处理 Sim-to-Real 的跨域迁移，展现出工业化的应用前景。
\end{enumerate}

\subsection{未来工作展望}
尽管 VLN-R1 在各项基准测试中表现优异，但在以下方面仍有进一步研究的空间：
\begin{enumerate}
    \item **多样化感知的融合**：目前系统主要依赖视觉与激光雷达（仅底层避障）。未来可考虑引入音频感知、触觉感知等多模态信息，提升智能体在极端动态环境下的自愈能力。
    \item **在线自进化能力**：目前模型主要依赖离线训练。探索如何使智能体在实际部署中通过在线学习（Online Learning）不断完善特定场景下的导航经验，是实现通用具身智能的关键。
    \item **长时程记忆管理**：当导航任务跨越多个楼层或超大尺度空间时，拓扑图的内存开销与检索效率将成为瓶颈。引入类似外部知识库或分层向量检索的技术将是未来的突破口。
\end{enumerate}
通过本文的工作，我们证明了将大语言模型的先进推理能力与具身控制相结合是实现新一代智能导航系统的可靠路径。
